\section{CLIP}
\textit{CLIP} (Contrastive Language-Image Pretraining) is a SOTA \textit{Vision-Language-Model (VLM)}, designed to learn a joint embedding space for images and text, 
enabling it to perform various tasks such as image classification, object detection, and zero-shot learning.

CLIP is composed of two main components:
\begin{itemize}
    \item An \textbf{image encoder}, typically a convolutional neural network (CNN) or a vision transformer (ViT), which processes images and produces feature representations.
    \item A \textbf{text encoder}, typically a transformer-based architecture, which processes text and produces feature representations.
\end{itemize}
The key innovation of CLIP is its use of a \textbf{contrastive learning approach} to train both encoders simultaneously. This approach allows
CLIP to learn a shared embedding space where semantically similar images and text are close together, while dissimilar pairs are far apart.

\paragraph{Before CLIP}
Before CLIP, the traditional learning approach for computer vision involved creating a training dataset of pairs of images and their corresponding labels.
This approach had a few limitations:
\begin{itemize}
    \item It required a large amount of labeled data, which can be expensive and time-consuming to collect.
    \item It was limited to the specific tasks for which it was trained, and could not easily generalize to new tasks or domains.
    \item To introduce new classes, the model had to be retrained on a new dataset, which was not efficient.
\end{itemize}
Since CLIP was trained on a large dataset of image-text pairs, it can perform zero-shot learning, meaning it can recognize and classify images based on textual descriptions without needing to be retrained on specific classes.

\subsection{Ideas that lead to CLIP}
In reality, CLIP didn't introduce anything revolutionary, but rather correctly combined existing ideas and innovations in the field of machine learning and computer vision. 
Some of the key ideas that led to the development of CLIP include:
\begin{itemize}
    \item \textbf{Webly supervised learning}: The idea of using large amounts of data from the web to train models, which has been a common practice in machine learning for many years.
    \item \textbf{Weakly supervised learning}: Using talked video data, which contains both visual and textual information, to learn a joint embedding space for images and text.
    \item \textbf{Going at scale}: The idea of training large models on massive datasets was already exploited in the field of Natural Language Processing (NLP) with models like GPT-3.
    CLIP applied this idea to the field of computer vision, demonstrating that scaling up the model and dataset can lead to significant improvements in performance.
\end{itemize}

\subsection{Dataset}
Previous works on this field had used datasets of limited size, such as MS COCO, which contains only 100K. 
A new dataset was needed to assess wether the proposed approach could scale up.

The dataset used to train CLIP, called \textit{WebImageText}, consists of 400 million image-text pairs collected from the internet.
It was created by scraping the web with more than 500k queries, built using a combination of common words, combinations of words, and random words. 
The queries were designed to be diverse and cover a wide range of topics and concepts, ensuring that the dataset is representative of the variety of images and text found on the internet.
To have an approximative class balance, each query includes up to 20k image-text pairs.

This dataset is significantly larger than previous datasets used for training vision-language models, but it is also noisier, 
as the data was minimally curated and quality of the image-text pairs is not guaranteed. 
This causes that the model reflects the biases and limitations of the data.

\subsection{Contrastive Pretraining}
As mentioned before, CLIP uses a contrastive learning approach to train both the image and text encoders simultaneously.
The training objective is to \textbf{maximize} the cosine similarity between the embeddings of the $N$ matching 
image-text pairs in a batch, while \textbf{minimizing} the cosine similarity of the $N^2 - N$ incorrect pairs.

The loss function used for training CLIP is a softmax-based contrastive loss, defined as:
\begin{equation}
    \mathcal{L} = -\frac{1}{2N} \sum_{i=1}^{N} \left(
    \underbrace{\log \frac{\exp(\operatorname{sim}(v_i, t_i) / \tau)}{\sum_{j=1}^{N} \exp(\operatorname{sim}(v_i, t_j) / \tau)}}_{\textit{Image} \to \textit{Text}}
    +
    \underbrace{\log \frac{\exp(\operatorname{sim}(t_i, v_i) / \tau)}{\sum_{j=1}^{N} \exp(\operatorname{sim}(t_i, v_j) / \tau)}}_{\textit{Text} \to \textit{Image}}\right),
\end{equation}

where $v_i$ and $t_i$ are the embeddings of the $i$-th image and text pair, $\text{sim}$ is the cosine similarity function, and $\tau$ is a temperature hyperparameter that controls the sharpness of the distribution.
This loss encourages the model to learn a shared embedding space where semantically similar images and text are
close together, while dissimilar pairs are far apart.

\paragraph{Temperature}: The temperature $\tau$ is a hyperparameter that controls the sharpness of the distribution in the contrastive loss.
A lower temperature makes the distribution sharper, which can lead to better performance but also increases the risk of overfitting. 
Conversely, a higher temperature makes the distribution softer over the possible classes, which can help to prevent overfitting but may lead to worse performance.

\paragraph{Training Phase}: Since the WebImageText dataset is very large, both the image and text encoders are trained from scratch, without using any pretraining on other datasets.
Simple data augmentation techniques are applied to the images.

\paragraph{Limitations of the training approach}
\label{CLIP:training-limitations}
This training approach has some limitations regarding the computational resources required. 
In particular, since we need to compute the cosine similarity between all pairs of images and text in a batch ($N^2$ pairs),
the size of the batch is limited by the available memory.

\subsection{Inference}
During inference, CLIP can be used for various tasks such as image classification, object detection, and zero-shot learning.
\begin{itemize}
    \item \textbf{Image classification}: Given an image, CLIP can be used to classify it by comparing its embedding to the embeddings of a set of candidate labels (text descriptions);
    \item \textbf{Text to Image retrieval}: Given a text query, CLIP can be used to retrieve relevant images by comparing the text embedding to the embeddings of a set of candidate images.
    First, the text query is encoded using the text encoder to obtain its embedding. Then, the cosine similarity between the text embedding and the embeddings of a set of candidate images is computed. 
    The images with the highest similarity scores are returned as the most relevant results.
    \item \textbf{Image to Image retrieval}: Given an image query, CLIP can be used to retrieve relevant images by comparing the image embedding to the embeddings of a set of candidate images.
    Similarly to text to image retrieval, the image query is encoded using the image encoder to obtain its embedding. Then, the cosine similarity between the image embedding and the embeddings of a set of candidate images is computed.
    The images with the highest similarity scores are returned as the most relevant results.
    \item \textbf{Zero-shot learning}: CLIP enables zero-shot learning, meaning it can recognize and classify images based on textual descriptions without needing to be retrained on specific classes.
\end{itemize}

\subsection{Fine Tuning}
Once the power of CLIP was demonstrated on zero-shot learning tasks, the next step was to explore how to further improve its performance on specific tasks.
A first approach was to fine-tune either the image encoder, the text encoder, or both on a specific dataset.
This was found to be ineffective, as the model had already learned a powerful shared embedding space, and fine-tuning on a specific dataset often led 
to overfitting and a decrease in performance on other tasks.

Other approaches based on prompt engineering were found to be more effective in improving the performance of CLIP on specific tasks, 
without the need for fine-tuning the model itself.
Here, we will discuss some of the most effective techniques for improving CLIP's performance through prompt engineering.

\subsubsection{Basic Prompt Engineering}:
Prompt engineering is the process of designing and optimizing the input prompts to improve a model's performance on specific tasks. 
Since CLIP learns a shared embedding space, the phrasing of the text prompts significantly impacts the results.

For example, since images are rarely described with a single word in the wild, using a template like \texttt{"a photo of a [class]"} 
often outperforms using the raw class name alone. On ImageNet, this simple modification yielded a gain of 1.3\% in accuracy.

\subsubsection{Prompt Ensembling}
Further improvements can be achieved by using multiple prompts for each class and averaging their embeddings. 
For the class "dog," one might ensemble prompts like \textit{"a photo of a dog,"} \textit{"a picture of a dog,"} and \textit{"a photo of a cute dog."} On ImageNet, 
this approach provided an additional 3.5\% gain in accuracy compared to a single prompt.

\subsubsection{Advanced Approaches}
Prompt engineering can still be difficult and time-consuming. Many times a slight change in the prompt can lead to a significant change in performance,
with no clear way to predict which prompts will work best for a given task.

To address this, more advanced approaches have been proposed.

\paragraph{CoOp}
CoOp (Context Optimization) is a method that learns continuous \textbf{prompt vectors} instead of using discrete text prompts.
These prompt vectors are then fed into the text encoder to generate the class embeddings.
Different choices can be made, such as the number of prompt vectors, their dimensionality, where they are inserted in the text encoder (e.g., before the class name, after the class name, or both), 
and whether they are shared across classes or specific to each class.

CoOp optimizes the prompt vectors using a small amount of labeled data from the target task, allowing it to adapt the prompts to better suit the specific task at hand.
This approach has been shown to significantly improve the performance of CLIP on various tasks, especially when the amount of labeled data is limited.

Even though the results are quite impressive, they are still not humanly interpretable, as the learned prompt vectors may not correspond to any meaningful words or phrases.
Also, the approach was shown to have poor generalization performance to new classes, as the learned prompt vectors may be overfitted to the specific classes seen during training.

\paragraph{CoCoOp}
CoCoOp (Contextualized Context Optimization) is an extension of CoOp that addresses the issue of generalization to new classes. 

CoCoOp learns a \textbf{meta-network} that generates some meta-prompts based on the output of the image encoder. These meta-tokens
are them summed with the learned prompt vectors to generate the final prompts that are fed into the text encoder.
This allows CoCoOp to adapt the prompts based on the specific image being processed, improving generalization to new classes and reducing overfitting to the classes seen during training.

The problem with this approach is that it is more computationally expensive than CoOp, as it requires an additional forward pass through the meta-network for each image,
during the training phase.

\paragraph{Visual Prompt Tuning}
All the approaches described so far have focused on optimizing the text prompts, while keeping the image encoder fixed.
Visual Prompt Tuning is a method that optimizes the input to the image encoder instead of the text encoder. 
Similarly to CoOp, it learns a set of continuous prompt vectors that are instead added to the input image before being fed into the image encoder.

\paragraph{Multimodal prompt Learning}
Sometimes, optimizing only the text prompts or only the image prompts leads to suboptimal performance, as ther's no general rule to determine which one is more effective for a given task.
Multimodal Prompt Learning is a method that optimizes both the text and image prompts simultaneously, allowing for a more holistic approach to prompt engineering.

Using a small neatwork (for example, a transformer), the method learns to generate both text and image prompts based on the input image and the target task.
The prompts are then fed into the respective encoders to generate the final embeddings, which are used for classification or other downstream tasks.

\subsection{Beyond CLIP - SigLIP}
To resolve the limitations of CLIP, presented in \autoref{CLIP:training-limitations}, a new method called SigLIP (Sigmoid Loss for Language-Image Pretraining) was proposed.
Instead of using a softmax-based contrastive loss, where each pair of image and text is compared to all other pairs in the batch, 
SigLIP uses a sigmoid-based loss that compares each pair independently, treating the problem as a binary classification.

For each pair of image and text, the model predicts a score that indicates whether the pair is a match or not (e.g., whether the image and text are semantically related).
Because each pair is treated independently, the batch size can be much larger without running into memory issues, allowing for more efficient training on large datasets.

